% from main.tex
%%%%%%%%%%  scan idx 400  |  folio 387  %%%%%%%%%%
\[
  \Projz\bigl(F_k(L^{\bullet})\bigr)\longrightarrow L^{\bullet},
\]
et c'est un
isomorphisme sur chaque composante puisque $L^{\bullet}$ a été supposé strictement parfait, d'où le ii).

\textbf{2.7.} Sur $\mathcal C$, le foncteur qui à un $B$-Module gradué associe sa $n$-ième composante est un foncteur exact à valeurs dans $\mathcal C'$; il définit donc un foncteur :
\begin{equation}\tag{2.7.1}
  \operatorname{Comp}_n:D^+(\mathcal C)\longrightarrow D^+(\mathcal C').
\end{equation}
On peut, pour tout $n$, considérer le foncteur $\mathcal C'\to\mathcal C$ défini par la tensorisation par $B(-n)$. L'homomorphisme (1.2.1) définit un morphisme de foncteurs sur $\mathcal C$, d'où un morphisme pour les foncteurs correspondants entre catégories dérivées :
\begin{equation}\tag{2.7.2}
  \operatorname{Comp}_n\otimes_{\mathcal O_S} B(-n)
  \longrightarrow \operatorname{Id}_{D^+(\mathcal C)}.
\end{equation}
Soit $M^{\bullet}$ un complexe $\mathcal C_0$-parfait. Nous allons montrer qu'il existe une suite de triangles distingués :
\[
\begin{array}{ccccccccccc}
& P_1^{\bullet} &&&& P_2^{\bullet} &&&& P_k^{\bullet} & \\
\swarrow && \nwarrow && \swarrow && \nwarrow && \swarrow && \nwarrow \\
N_1^{\bullet} & \longrightarrow & M^{\bullet}, & \qquad &
N_2^{\bullet} & \longrightarrow & P_1^{\bullet}, & \cdots, &
N_k^{\bullet} & \longrightarrow & P_{k-1}^{\bullet},
\end{array}
\]
où $P_k^{\bullet}$ est acyclique, et les $N_i^{\bullet}$ du type \(K_i^{\bullet}\otimes_{\mathcal O_S}B(d)\), les $K_i^{\bullet}$ étant des complexes parfaits au sens ordinaire sur $S$.

Puisque $M^{\bullet}$ est $\mathcal C_0$-parfait, et que $S$ est quasi-compact, il existe un recouvrement fini de $S$, et un entier $d_0$ tel que sur tout ouvert de ce recouvrement, $M^{\bullet}$ soit isomorphe à un complexe borné à termes localement libres de type fini sur $B$, soit $M_U'{}^{\bullet}$, tel que $\operatorname{Comp}_d(M_U'{}^{\bullet})=0$ pour $d<d_0$.

